\documentclass[10pt,a4paper]{amsart}
\usepackage[margin=19mm]{geometry}
\usepackage{amsmath,amssymb,mathtools}
\usepackage[scr=rsfs,cal=euler]{mathalfa}
\usepackage{xfrac}
\usepackage[unicode,hidelinks,bookmarks=false]{hyperref}
\newcommand{\ie}{i.e.\ }
\newcommand{\cf}{cf.\ }
\newcommand{\resp}{respectively\ }
\newcommand{\Subsection}{\subsection}
\providecommand{\nobreakdash}{\nobreak\hbox{-}}
\setlength{\emergencystretch}{2em}
\multlinegap0pt
\usepackage{amssymb}
\usepackage{mathtools}
\usepackage[shortlabels]{enumitem}
\newtheorem{proposition}{Proposition}
\newcommand{\Q}{\ensuremath{\mathbb{Q}}}
\newcommand{\Z}{\ensuremath{\mathbb{Z}}}
\newcommand{\clq}{C}
\newcommand{\mapsfrom}{\mathrel{\reflectbox{\ensuremath{\mapsto}}}}
\title{On exceptional cliques in matrix rings}
\author{Milan Boutros, Ignacio Cascudo, Ronald Cramer, Dani\"el van Gent, Chaoping Xing}
\date{Source: arXiv:2608.24586v1 (CC BY 4.0)}
\begin{document}
\maketitle

\noindent\emph{Source and licence.} On exceptional cliques in matrix rings. Version arXiv:2608.24586v1 (CC BY 4.0), distributed by arXiv under the Creative Commons Attribution 4.0 International licence, http://creativecommons.org/licenses/by/4.0/, as recorded in the arXiv licence metadata for this exact version and shown on the abstract page https://arxiv.org/abs/2608.24586v1. Full licence notice: \texttt{licenses/arxiv-2608.24586-CC-BY-4.0.txt}.\par\smallskip

\noindent\textit{Editorial scope.} Counted unit: the article's proposition prop:dim2\_conj (source0.tex lines 1561--1563). The article's complete original proof of it is delivered with it (source0.tex lines 1564--1659, one \texttt{proof} environment of 96 lines). No mathematics is written, repaired or re-derived: the statement and the proof are the article's own lines, in the article's own notation, and no retained line is substituted or re-flowed.  No further source-local context is needed: every cross-reference of every retained line resolves inside the two delivered spans.  Nothing was omitted from any retained span, so the package carries no omission record.  No retained line contains a citation, so the wrapper carries no bibliography and no bibliography entry.  No retained line contains a figure environment, an image-inclusion command or drawing code, so no image is delivered and the package declares no asset dependency.  Editorial formatting only, in addition: an AMS \texttt{amsart} preamble that restates the article's own declarations and macros used by the retained lines, namely the environments \texttt{proposition} on separate counters, together with the standard packages those lines need.  The retained environments print in this document as Proposition 1 in order of appearance, whereas the article prints the same items with the numbering of its own full document.  The complete native source of the article as distributed in the arXiv e-print for this exact version is bundled unchanged as \texttt{sources/208492/source0.tex}.
	\begin{proposition}\label{prop:dim2_conj}
		Each exceptional cliques of size \(4\) in \(\textup{Mat}_2(\Z)\) is equivalent to either \(\{0,1,\phi,\phi+1\}\) or \(\{0,1,\zeta,1-\phi\}\).
	\end{proposition}

	\begin{proof}[Sketch]
		The two cliques are non-equivalent, since the first is commutative while the second is not. 
		Let \(u=\big(\begin{smallmatrix}0&1\\s&1\end{smallmatrix}\big)\) for \(s\in\{1,-1\}\).
		Fix \((d_0,d_1,d_u)\in\{-1,1\}^3\).
		We will show that we have bijections 
		\begin{align*} 
			\delta&=d_1-d_0-1 \\
			\varepsilon &= d_u-d_0+s \\
			x_0&=\delta-2\varepsilon, \\
			y_0&=\delta, \\
			D&=1+4s, \quad\text{and}\\
			N&=x_0^2-D(y_0^2-4d_0).
		\end{align*}
		We have a bijection
		\begin{align*}
			\{v\in\textup{Mat}_2(\Z) \,|\, (\forall\, t\in\{0,1,u\})\ \det(v-t)=d_t \} &\leftrightarrow \{(a,b)\in\Z^2 \,|\, a^2+ab-sb^2+\delta a+\epsilon b+d_0=0 \}, \\
			%\big(\begin{smallmatrix}a&b\\\clq&d\end{smallmatrix}\big) &\mapsto (a,b), \\
			\big(\begin{smallmatrix}a&b\\a-sb+\varepsilon&-a-\delta\end{smallmatrix}\big) &\mapsfrom (a,b)
		\end{align*}
		and a bijection
		\begin{align*}
			\{(a,b)\in\Z^2 \,|\, a^2+ab-sb^2+\delta a+\epsilon b+d_0=0 \} &\leftrightarrow \{  z \in \Z[\sqrt{D}] \,|\, N(z)=N,\, z \equiv x_0\ (\textup{mod }\sqrt{D}) \}, \\
			(a,b) &\mapsto (Db+x_0) + (2a+b+y_0)\sqrt{D}.
		\end{align*}
		
		\underline{Case \(u=\zeta\):} If \(s=-1\), then \(D=-3<0\) and the norm equation \(N(z)=N\) has at most \(\#\Z[\sqrt{D}]^*=2\) solutions. Hence there are at most \(12\) possible \(v\) such that \(\{0,1,u,v\}\) is an exceptional clique. In fact, we have the following \(12\) solutions:
		\begin{align*}
			\left(\begin{smallmatrix}
				-1 & \phantom{-}1 \\
				-3 & \phantom{-}2
			\end{smallmatrix}\right),
			\left(\begin{smallmatrix}
				-1 & \phantom{-}3 \\
				-1 & \phantom{-}2
			\end{smallmatrix}\right),
			\left(\begin{smallmatrix}
				\phantom{-}1 & \phantom{-}1 \\
				-1 & \phantom{-}0
			\end{smallmatrix}\right),
			\left(\begin{smallmatrix}
				-1 & \phantom{-}1 \\
				-1 & \phantom{-}2
			\end{smallmatrix}\right),
			\left(\begin{smallmatrix}
				\phantom{-}1 & -1 \\
				-1 & \phantom{-}0
			\end{smallmatrix}\right),
			\left(\begin{smallmatrix}
				\phantom{-}1 & \phantom{-}1 \\
				\phantom{-}1 & \phantom{-}0
			\end{smallmatrix}\right), \\
			\left(\begin{smallmatrix}
				-1 & -1 \\
				-1 & \phantom{-}0
			\end{smallmatrix}\right),
			\left(\begin{smallmatrix}
				-1 & \phantom{-}1 \\
				\phantom{-}1 & \phantom{-}0
			\end{smallmatrix}\right),
			\left(\begin{smallmatrix}
				\phantom{-}1 & -1 \\
				\phantom{-}1 & -2
			\end{smallmatrix}\right),
			\left(\begin{smallmatrix}
				\phantom{-}1 & -1 \\
				-1 & \phantom{-}2
			\end{smallmatrix}\right),
			\left(\begin{smallmatrix}
				\phantom{-}1 & \phantom{-}1 \\
				\phantom{-}1 & \phantom{-}2
			\end{smallmatrix}\right),
			\left(\begin{smallmatrix}
				\phantom{-}3 & -1 \\
				\phantom{-}1 & \phantom{-}0
			\end{smallmatrix}\right).
		\end{align*}
		The transformations \(x\mapsto u x u^{-1}\), \(x\mapsto (1-x)(1-u)^{-1}\) and \(x \mapsto u x^{-1}\) generate some ERROR group \(G\) which preserves the exceptional clique \(\{0,1,u\}\) and transitively permutes the 12 solutions. Hence \(\{0,1,u,v\}\) is equivalent to \(\{0,1,\zeta,\phi+1\}\).
		
		\underline{Case \(u=\phi\):} Suppose \(s=1\). Then \(D=5\), and \(\eta=\tfrac{1}{2}(1+\sqrt{5})\) is a fundamental unit of the maximal order \(\mathcal{O}\) of \(\Q(\sqrt{5})\). With \(U=\{z\in\mathcal{O}^*\,|\, z \equiv 1\ (\textup{mod } \sqrt{5})\} = \langle-\eta^2\rangle\), the solutions \(z\in\Z[\sqrt{5}]\) to \(N(z)=z\) and \(z\equiv x_0\ (\text{mod \(\sqrt{5}\)})\) are as follows:
		\begin{align*}
			\begin{array}{ccc|ll|l}
				d_0 & d_1 & d_u & x_0 & N & z \\ \hline
				-1 & -1 & -1 & -3 & -2^4 & (2+2\sqrt{5})\cdot U\\
				-1 & -1 & +1 & -7 & +2^3\cdot 3 & \text{none, \(3\) is inert}\\
				-1 & +1 & -1 & -1 & -2^3 \cdot 3 & \text{none, \(3\) is inert}\\
				-1 & +1 & +1 & -5 & \phantom{+}0 & 0\\
				+1 & -1 & -1 & -1 & -2^3 \cdot 3 & \text{none, \(3\) is inert}\\
				+1 & -1 & +1 & -5 & \phantom{+}0 & 0\\
				+1 & +1 & -1 & +1 & +2^4 & -4\cdot U \\
				+1 & +1 & +1 & -3 & +2^3\cdot 3 & \text{none, \(3\) is inert}\\
			\end{array}
		\end{align*}
		The rows where \(z=0\) correspond to the solutions \(v=\phi-1\) and \(v=\phi+1\) respectively, which gives a clique equivalent to \(\{0,1,\phi,\phi+1\}\).
		For the two rows with infinitely many solutions, one verifies that multiplying a solution by \(-\eta^2\) corresponds to conjugation by \(\phi-1\) in \(\textup{Mat}_2(\Z)\).
		Hence up to equivalence the cases \(v=\left(\begin{smallmatrix}1 & 1 \\1 & 0\end{smallmatrix}\right)\) and \(v=\left(\begin{smallmatrix}1 & -1 \\1 & 0\end{smallmatrix}\right)\) remain. The latter matrix is conjugate to \(\zeta\), and thus the clique \(\{0,1,\phi,v\}\) is equivalent to \(\{0,1,\zeta,\phi+1\}\) by the previous case. Finally, note that also \(\phi^{-1}\cdot \left(\begin{smallmatrix}1 & 1 \\1 & 0\end{smallmatrix}\right)\) is conjugate to \(\zeta\), so \(\{0,1,\phi,v\}\) too is equivalent to \(\{0,1,\zeta,\phi+1\}\).
	\end{proof}

\end{document}
